\begin{theorem}[S-33: Odd ramification obstructs the balanced character in every rank]
\label{thm:semireg-odd-ramification}
Let \((X,\theta)\) be a polarized complex abelian fourfold, with polarization volume \[
 D=\int_X\frac{\theta^4}{24}\in\mathbf Z_{>0}.
\] Let \(F=\mathbf Q(\sqrt d)\) be a real quadratic field, where \(d>1\) is squarefree. Assume its rational action is fixed by the Rosati involution and its two complex eigenspaces each have dimension two. Write \[
 \theta=\theta_++\theta_-,\qquad
 C=\frac{\theta_+^2+\theta_-^2}{2},\qquad
 \nu=\frac{\theta^4}{24}
      =\frac{\theta_+^2\theta_-^2}{4},\qquad \int_X\nu=D.
 \tag{S33.1}
\] For \(r\in F\), let \(\theta_r\) denote the rational alternating form \((x,y)\mapsto\theta(rx,y)\), equivalently \(\sigma_+(r)\theta_++\sigma_-(r)\theta_-\). Fix an odd prime \(p\) ramified in \(F\), so \(p\mid d\), and assume \[
 p\nmid D,\qquad
 (\mathcal O_F\otimes\mathbf Z_p)H_1(X,\mathbf Z_p)
                  \subseteq H_1(X,\mathbf Z_p).
 \tag{S33.2}
\] Thus only maximal-order action at this prime is required.

There is \textbf{no} perfect object \(\mathcal E\) of positive rank \(N\), with \(\ell=c_1(\mathcal E)\in F\theta\), and no \(u\in\mathbf Q^\times\) with \(v_p(u)=0\), satisfying \[
 e^{-\ell/N}\operatorname{ch}(\mathcal E)
       =N-uC+\frac{u^2}{N}\nu.
 \tag{S33.3}
\] The assertion allows an arbitrary fractional normalizing class \(-\ell/N\) and does not require \(\mathcal E\) to be a sheaf, simple or gluable. For the balanced character \(Nv_q\), the scalar here is \(u=Nq\).

In particular, suppose \(\operatorname{NS}(X)_{\mathbf Q}=F\theta\) and maximal-order action is integral at every ramified prime. For a simple gluable fully semiregular lift with an effective finite coherent target action \(G\) as in S-22, every odd ramified prime must divide \(D|G|\). Combining with S-32's even-ramification obstruction gives the necessary condition \[
 \operatorname{rad}(|\operatorname{Disc}(F)|)\mid D|G|.
 \tag{S33.4}
\] In the nonequivariant case, all ramified primes must divide \(D\). In particular, there is no such positive-rank realization for a principal polarization and a genuine real quadratic field. This excludes the proposed normalized character under the stated local maximal-order and determinant hypotheses; it does not assert nonexistence for arbitrary derived characters or arbitrary polarizations/actions.
\end{theorem}
\begin{proof}
\paragraph{1. The integral endomorphism corresponding to the determinant}

Since \(\ell\in F\theta\), write \(\ell=\theta_r\) for \(r\in F\). Because \(p\nmid D\), the integral alternating form \(\theta\) is perfect over \(\mathbf Z_p\): its determinant is \(D^2\) up to a unit in any integral basis. Integrality of \(\ell\) therefore implies \(rH_1(X,\mathbf Z_p)\subseteq H_1(X,\mathbf Z_p)\). Put \(K=F\otimes\mathbf Q_p\), \(O=\mathcal O_F\otimes\mathbf Z_p\) and \(L=H_1(X,\mathbf Z_p)\). The ring \(O\) is the valuation ring of a ramified quadratic extension. The \(O\)-module \(L\) is torsion-free of rank four, hence free. The scalar \(r\) preserves \(O^4\), and consequently \(r\in O\).

Set \(s=\sqrt d\). Because \(p\) is odd, \(O=\mathbf Z_p[s]\); the different is the principal ideal \((s)\), and \(s\) is a uniformizer up to a unit. Normalize its valuation by \(v_\pi(s)=1\), so \(v_\pi(p)=2\). We can write \[
 r=a+bs,\qquad a,b\in\mathbf Z_p.
 \tag{S33.5}
\]

\paragraph{2. The self-dual trace-symplectic normal form}

For completeness, the precise normal form used in the coefficient calculation is justified here. The Rosati condition says \(\theta(ax,y)=\theta(x,ay)\) for \(a\in K\). Nondegeneracy of the field trace determines a unique alternating \(K\)-bilinear form \(\beta\) on \(L\otimes\mathbf Q_p\) by \[
 \operatorname{Tr}_{K/\mathbf Q_p}(a\beta(x,y))=\theta(ax,y)
 \quad(a\in K).
\] The balance identity proves bilinearity in both variables. Alternation follows from \(\theta(ax,x)=0\), using the balance identity and skew symmetry in characteristic zero.

For \(x,y\in L\), integrality after multiplying \(x\) by every \(a\in O\) gives \(\beta(x,y)\in O^\vee=s^{-1}O\), where \(O^\vee\) is the trace-dual lattice. The trace pairing identifies \(\operatorname{Hom}_O(L,s^{-1}O)\) with \(\operatorname{Hom}_{\mathbf Z_p}(L,\mathbf Z_p)\). Since the local form \(\theta\) is perfect, \(\beta\) is a perfect \(s^{-1}O\)-valued alternating form. Thus \(\psi=s\beta\) is a unimodular alternating \(O\)-bilinear form.

Choose a symplectic \(O\)-basis \((e_1,f_1,e_2,f_2)\) for \(\psi\). One obtains it by choosing a primitive vector, a second vector pairing to a unit, rescaling that vector to make the pairing one, and splitting off the orthogonal hyperbolic plane; the residual rank-two summand is handled in the same way. No division by two is involved. In this basis \[
 L=U_1\oplus U_2,
 \quad U_i=Oe_i\oplus Of_i,
 \quad \psi(e_i,f_j)=\delta_{ij},
 \quad \theta=\operatorname{Tr}_{K/\mathbf Q_p}(\psi/s).
 \tag{S33.6}
\] In particular each \(U_i\) is a unimodular symplectic \(\mathbf Z_p\)-lattice of rank four for the restricted trace form.

\paragraph{3. The mixed second-Chern-class coefficients}

The degree-four part of (S33.3) gives \[
 \operatorname{ch}_2(\mathcal E)=\frac{\ell^2}{2N}-uC.
\] Since \(\ell^2/2\) is integral on a torus and \(c_2(\mathcal E)\) is integral, the identity \(\operatorname{ch}_2=\ell^2/2-c_2\) shows \[
 uC-\frac{\ell^2}{2N}\in H^4(X,\mathbf Z_p).
 \tag{S33.7}
\] For \(x,y\in O\), evaluate this four-form on \((xe_1,f_1,ye_2,f_2)\). The two embedding components of the trace form give \[
 C(xe_1,f_1,ye_2,f_2)=\operatorname{Tr}_{K/\mathbf Q_p}(xy/d),
\] while the orthogonality of \(U_1,U_2\) gives \[
 (\ell^2/2)(xe_1,f_1,ye_2,f_2)
 =\operatorname{Tr}(rx/s)\operatorname{Tr}(ry/s).
 \tag{S33.8}
\] Indeed each embedding component on one \(O\)-hyperbolic plane has rank two, so its square vanishes there; the mixed term in \(\theta_\sigma^2/2\) is the product of its two plane restrictions, with coefficient one.

In the basis \((1,s)\) of \(O\), the trace Gram matrix in (S33.8) for \(C\) is \(\operatorname{diag}(2/d,2)\). The corresponding linear functional for \(\ell\) has values \(2b,2a\). Therefore (S33.7) implies \[
 \frac{b^2}{N}-\frac{u}{2d}\in\mathbf Z_p,
 \qquad \frac{a^2}{N}\in\mathbf Z_p,
 \qquad \frac{ab}{N}\in\mathbf Z_p.
 \tag{S33.9}
\] Here division by four is legitimate because \(p\) is odd. These displayed factors are tied to the convention \(\theta=\operatorname{Tr}(\psi/s)\); no \(1/(2s)\) normalization has been substituted.

Let \(m=v_p(N)\) and \(k=v_p(b)\). The first congruence forces \(b\ne0\) and \[
 m=2k+1>0,
 \qquad v_p(a)\ge k+1,
 \qquad v_\pi(r)=2k+1=m.
 \tag{S33.10}
\] The summands \(a\) and \(bs\) have distinct \(\pi\)-valuations, at least \(2k+2\) and exactly \(2k+1\), respectively, proving the last equality. Moreover \[
 \frac{r^2}{N}
 =\frac{a^2}{N}+d\frac{b^2}{N}
             +2\frac{ab}{N}s
 \equiv\frac u2\pmod\pi.
 \tag{S33.11}
\] The first summand has positive \(p\)-valuation, the second has the asserted residue by (S33.9), and the third is divisible by \(\pi\). In particular \(u-r^2/N\) is a unit of \(O\).

\paragraph{4. The integral Euler characteristic gives the contradiction}

The tangent bundle of an abelian variety is trivial, so Riemann--Roch gives \(\chi(\mathcal E)=\int_X\operatorname{ch}_4(\mathcal E)\in\mathbf Z\). The character (S33.3) yields \[
 \operatorname{ch}_4(\mathcal E)
 =\frac{\ell^4}{24N^3}
       -\frac{u\ell^2C}{2N^2}+\frac{u^2}{N}\nu.
\] Using (S33.1) and writing \(r_\pm=\sigma_\pm(r)\), one computes \[
 \int_X\frac{\ell^4}{24}=D(r_+r_-)^2,
 \qquad
 \int_X\ell^2C=2D(r_+^2+r_-^2).
\] It follows that \[
 \chi(\mathcal E)
 =D\frac{\operatorname{Nm}(r)^2-Nu\operatorname{Tr}(r^2)+N^2u^2}{N^3}
 =D\frac{\operatorname{Nm}_{F/\mathbf Q}(Nu-r^2)}{N^3}.
 \tag{S33.12}
\] By (S33.11), \(Nu-r^2=N(u-r^2/N)\) has \(\pi\)-valuation \(2m\). Its local field norm has \(p\)-valuation \(2m\), since this ramified quadratic extension has residue degree one. Therefore \[
 v_p(\chi(\mathcal E))=2m-3m=-m<0,
 \tag{S33.13}
\] contradicting the integer Euler characteristic. This proves the main assertion.

For the semiregularity corollary, S-22 gives \(4D(Nq)^2/|G|=k\in\{2,4\}\). If an odd ramified prime does not divide \(D|G|\), then \(u=Nq\) is a \(p\)-adic unit. The degree-two character equation forces \(B=-\ell/N\), so the proposed normalized character is exactly (S33.3). The Néron--Severi hypothesis gives \(\ell\in F\theta\). The contradiction already proved shows that every odd ramified prime divides \(D|G|\). At \(2\), when \(2\nmid D|G|\), the value \(k=2\) would give the rational square \(u^2\) odd \(2\)-adic valuation \(-1\), which is impossible; \(k=4\) gives a \(2\)-adic unit, to which S-32 applies. This proves the combined condition (S33.4).
\end{proof}
\begin{remark}[An independent third-character certificate]
There is also a local contradiction before using the top-degree Euler calculation. This records the exact coefficient and its sign.

Let \(\theta_{\sigma,i}\) be the embedding component of \(\theta|_{U_i}\). Put \[
 \eta_1=\theta_{+,1}\theta_{-,1}
       =\theta|_{U_1}^2/2.
\] This is a primitive \(\mathbf Z_p\)-volume form on \(U_1\). In the particular ordered basis \((e_1,se_1,f_1,sf_1)\) its coefficient is \(4\) up to sign, which is harmless at odd \(p\). The degree \((4,2)\) component for the direct sum \(U_1\oplus U_2\) satisfies \[
 \begin{aligned}
 (\ell^3/6)_{(4,2)}
    &=\operatorname{Nm}(r)\eta_1\ell|_{U_2},\\
 (\ell C)_{(4,2)}
    &=\eta_1\theta_{\bar r}|_{U_2}.
 \end{aligned}
\] The second formula follows by writing \(C=\theta_{+,1}\theta_{+,2}+\theta_{-,1}\theta_{-,2}\): wedging its first summand with \(\ell|_{U_1}\) picks out \(r_-\), and wedging its second picks out \(r_+\). Thus \[
 \bigl(\operatorname{ch}_3(\mathcal E)\bigr)_{(4,2)}
 =\eta_1\theta_S|_{U_2},
 \qquad
 S=\frac{\operatorname{Nm}(r)r}{N^2}-\frac{u\bar r}{N}
   =\frac{\bar r(r^2-Nu)}{N^2}.
 \tag{S33.14}
\] For any integral degree-two form \(\ell\) on a torus, all divided powers \(\ell^j/j!\) are integral: in an exterior basis the coefficient of each ordered \(2j\)-monomial is a Pfaffian with integral entries. Hence \[
 \operatorname{ch}_3(\mathcal E)
 =\frac{\ell^3}{6}-\frac{\ell c_2(\mathcal E)}2
                         +\frac{c_3(\mathcal E)}2
 \in H^6(X,\mathbf Z_p)
\] for every odd \(p\), including \(p=3\). Since \(\eta_1\) is primitive and the restricted local trace form is perfect, (S33.14) would imply \(S\in O\). But (S33.10)--(S33.11) imply \[
 v_\pi(S)=v_\pi(\bar r)-v_\pi(N)
                  +v_\pi(r^2/N-u)=m-2m+0=-m<0.
\] This independently contradicts integrality. The argument uses only the usual integral Chern classes of a perfect object and the exterior algebra of a torus; no general integral topological Chern-character theorem is needed.
\end{remark}
